\chapter{Integración económica}

\section{Globalización}

En las últimas décadas, ha aumentado la interdependencia entre las economías del mundo. Esta interdependencia se manifiesta como movimientos internacionales de bienes y servicios, trabajo, empresas, fondos de inversión y tecnología \parencite[p.~1]{Carbaugh2009}. Muestra de ello es el crecimiento de la importancia relativa de las exportaciones e importaciones como proporción del \gls{pib} mundial, que pasó del 25\% en 1970 al 56.5\% en 2021. Por su parte, la proporción de salidas de \gls{ied} pasó de 0.6\% a 2.2\% en el mismo periodo \parencite{BancoMundial}.

El proceso de globalización\footnote{El concepto de globalización es polémico en su definición y periodización. Pueden encontrarse otras perspectivas sobre este concepto en \textcite{Ianni2004, Estay2016}.} ``consiste en la creciente comunicación e interdependencia entre distintos países del mundo, unificando sus mercados, sociedades y culturas, a través de una serie de transformaciones sociales, económicas y políticas'' \parencite[p.~222]{Perez2012}. En el plano económico, la globalización implica una mayor integración de los mercados de productos y servicios entre los países, por medio del comercio, la migración laboral y los flujos de capital \parencite[p.~2]{Carbaugh2009}.

Dentro de las causas de la globalización en la economía, podemos encontrar el avance en las tecnologías de la información y la comunicación, y su efecto en la reducción de los costos de transporte y comercialización; la apertura paulatina de las fronteras económicas a partir de la segunda guerra mundial; el crecimiento de los movimientos de capital, particularmente entre los países de la \gls{ocde} desde la década de los noventa; así como, la liberalización interior de muchas economías nacionales \parencite[p.~343--345]{Requijo2011}.

La globalización no es un fenómeno nuevo, han existido olas anteriores de globalización; por ejemplo, entre los periodos de 1870 a 1914 y entre 1945 y 1980. En el último caso, se caracterizó por la especialización de los países ricos en la industria manufacturera y el aumento de la productividad gracias a las economías de aglomeración \parencite[p.~4--5]{Carbaugh2009}. La principal diferencia es que la globalización actual es mucho más profunda y amplia que hace un siglo \parencite{Bordo1999}. Los principales rasgos que diferencian el proceso actual de globalización son: el aumento rápido de los flujos comerciales, especialmente de flujos transfronterizos de capital; la relevancia creciente de las empresas multinacionales; la deslocalización productiva; las crisis financieras recurrentes; la competencia mundial para la mayoría de los productos, tanto financieros como no financieros; y la escasa funcionalidad de las políticas nacionales \parencite[p.~346]{Requijo2011}.
%Bordo p.~3--4 menciona textualmente, la frase

\section{Integración económica}

Junto con la globalización, aparecen procesos de integración económica que refieren a diferentes tipos de acuerdos mediante los cuales los países eliminan de forma progresiva las barreras al comercio, ya sea bilateral o multilateral \parencite[p.~215]{Perez2012}. El grado de integración depende de los compromisos que adquieren los países a través de sus acuerdos comerciales o aduaneros \parencite[p.~29--31]{Witker2011}, los cuales pueden ser:
\begin{itemize}
	\item Acuerdo preferencial: Los países miembros se conceden mutuamente preferencias, como reducción de aranceles, en sectores específicos.
	\item Tratado de libre comercio: Los miembros eliminan los aranceles y otros obstáculos comerciales entre sí; aunque mantienen aranceles nacionales frente a terceros países.
	\item Unión aduanera: Los integrantes eliminan todos los obstáculos a la libre circulación de bienes entre sí, y establecen un arancel común y otras fórmulas protectoras frente a terceros países.
	\item Mercado común: Consiste en una unión aduanera  con libre movilidad de los factores de producción (capital y mano de obra). En algunos casos, también pueden existir políticas públicas comunes.
	\item Mercado único: Es un perfeccionamiento del mercado común en el que se eliminan las fronteras físicas, técnicas y fiscales en la circulación de bienes y servicios.
	\item Unión económica: Los miembros establecen políticas macroeconómicas coordinadas.
	\item Unión monetaria: Corresponde a una unión económica en la que se establece una única moneda común, que sustituye a todas las monedas nacionales de los miembros.
	\item Unión económica plena: Es el máximo nivel de integración económica de los países miembros que tiene como consecuencia la unión política de sus integrantes \parencite[p.~38--39]{Requijo2011}.
\end{itemize}
En \textcite{Tamames2012} puede encontrarse una revisión histórica de diferentes ejemplos de procesos de integración económica en el mundo.

La base de datos sobre \gls{rta} de la \gls{omc} contabiliza 378 acuerdos en vigor. La tabla \ref{Tabla_AcuerdosComerciales} muestra su distribución por tipo de acuerdo.

\begin{table}[h]
	\centering
\begin{threeparttable}
	\caption{Acuerdos comerciales regionales en vigor \parencite{WTO_RTA}}
	\label{Tabla_AcuerdosComerciales}
	\begin{tabular}{llc}
		\hline
		\textbf{Tipo de acuerdo\tnote{a}} & \textbf{Cobertura} & \textbf{Cantidad} \\
		\hline
		CU & Bienes & 11 \\
		CU y EIA & Bienes y servicios & 5 \\
		EIA & Servicios & 3 \\
		FTA & Bienes & 133 \\
		FTA y EIA & Bienes y servicios & 199 \\
		PSA & Bienes & 26 \\
		PSA y EIA & Bienes y servicios & 1 \\
		\hline
	\end{tabular}
	
	\begin{tablenotes}
		\item[a] \footnotesize{\gls{cu}, \gls{eia}, \gls{fta2}, \gls{psa}.}
	\end{tablenotes}
	
\end{threeparttable}
\end{table}

% que pueden ser zona de libre comercio, unión aduanera, mercado común o unión económica \parencite[p.~29--31]{Witker2011}.%\footnote{Las características de cada nivel de integración pueden encontrarse en \parencite{Requijo2011} 
 
% Explicar cada concepto y ejemplos

Existe cierta tensión entre la integración regional y la liberalización de los mercados a escala mundial, en tanto que las agrupaciones regionales afectan los mecanismos de implantación de un sistema internacional abierto, eficiente y multilateral, representado por la \gls{omc} \parencite[p.~222]{Perez2012}. \textcite{Krugman1992} mostró mediante un modelo sencillo que el comportamiento óptimo no cooperativo de una economía mundial conformada por bloques económicos es incrementar los aranceles externos. Bajo su modelo simétrico, la formación de bloques comerciales, con libre comercio interno y aranceles comunes externos, conduce al aumento de aranceles y a la reducción del comercio entre bloques.

\section{Cadenas Globales de Valor}

Una característica de la globalización es la deslocalización y fragmentación de la producción de bienes y servicios a nivel mundial. Esta nueva fase de la organización productiva es denominada formación de \acrfull{cgv}. Esto representó un cambio con las estrategias de integración vertical, ya que las grandes empresas multinacionales optaban por deslocalizar y externalizar fases de sus procesos a otros países \parencite[p.~63]{Perez2019}. Sin embargo, la fragmentación internacional de la producción no se distribuyó de manera uniforme en el mundo, sino que muestra un patrón de concentración regional acentuado \parencite[p.~30--31]{Porta2017}.

De acuerdo con el \gls{gbm},\footnote{El \acrlong{gbm} está conformado por el \gls{birf}, la \gls{aif}, la \gls{cfi}, el \gls{omgi} y el \gls{ciadi}, véase \textcite{GrupoBancoMundial_QUienesSomos}.} la formación de \glspl{cgv} ha permitido a los países en desarrollo especializarse y crecer, sin tener que crear industrias completas desde cero. Las \glspl{cgv} han favorecido que las empresas de los países en desarrollo mejoren su productividad y han sido un factor de ayuda para su transición de las exportaciones de productos básicos a la producción de manufacturas básicas \parencite{GrupoBancoMundial2019}.

Las \glspl{cgv} pueden definirse como ``un conjunto de eslabones consecutivos de actividades económicas en diferentes países o regiones, que interactúan en una dinámica continua'' \parencite[p.~2]{Ferrando2013}. De esta manera, el proceso de producción se divide entre distintos países, donde sus empresas se especializan en una tarea específica sin elaborar todo el producto \parencite[p.~3]{BancoMundial2020}. El proceso de formación de \acrshortpl{cgv} es la más reciente manifestación de la integración económica en el mundo: anteriormente, el comercio internacional era principalmente de bienes y servicios finales, pero la reducción de los costos de transporte y la liberalización económica propiciaron el crecimiento del comercio de bienes de consumo y la fragmentación del proceso productivo a escala global \parencite{BancoMexico2016}.

El estudio de las \acrshortpl{cgv} ha tomado gran relevancia en las agendas de investigación de organismos internacionales como la \gls{ocde}, la \gls{omc} y \gls{unctad}, e inclusive en los principales foros económicos internacionales \parencite[p.~3]{Dalle2013}. El creciente interés de estas organizaciones por las \acrshortpl{cgv} impulsó la generación de nuevas bases de datos que permiten analizar la influencia de estas sobre la economía mundial; por ejemplo, mediante la creación de matrices insumo--producto a nivel mundial \parencite[p.~31]{Porta2017}. 

% 2023 Global Value Chain Development report
% 2020 Banco Mundial_Informe sobre el desarrollo mundial 2020

Cuantificar el comercio internacional en el marco de las \acrshortpl{cgv} requiere nuevas formas de medir el intercambio de bienes y servicios; véase \textcite{Jones2014, Timmer2015}. Las transacciones de comercio exterior en valores brutos pueden generar problemas de doble contabilidad, debido a que los bienes usados como materia prima o insumos intermedios pueden cruzar varias veces las fronteras \parencite[p.~6]{Dalle2013}. Para contabilizar la influencia de las \acrshortpl{cgv} en la generación de valor agregado en distintos países se utilizan las matrices insumo--producto mundiales, que combinan y armonizan la información de las matrices de insumo--producto nacionales con los flujos de comercio bilateral de bienes y servicios \parencite{BancoMexico2016, Degain2016}.%Degain p.~19


\section{Matrices Insumo--Producto}

Las matrices de insumo--producto (\acrshort{mip}) son una herramienta metodológica para describir las interrelaciones productivas en una economía durante un periodo de tiempo, generalmente un año \parencite[p.~13]{Durán2021}. Este marco analítico fue desarrollado por Wassily Leontief a finales de la década de 1930, por lo cual fue reconocido con el premio Nobel de Economía en 1973. En su forma básica, el modelo de insumo--producto consiste en un sistema de ecuaciones lineales que describen la distribución del producto de cada industria en la economía \parencite[p.~1]{Miller2012}. En un artículo de 1936, Leontief planteó el modelo insumo--producto como un sistema de contabilidad de doble entrada, que registra en su lado acreedor las salidas de bienes y servicios de una industria, y en su lado deudor las adquisiciones de bienes o servicios. Su objetivo era construir una \textit{Tableau Économique} de los Estados Unidos en 1919 que describiera la interdependencia productiva del país \parencite[p.~105]{Leontief1936}. Al año siguiente, desarrolló una base teórica de este modelo a partir de la teoría económica del equilibrio general, mediante un conjunto de sistemas de ecuaciones lineales. Los primeros tres conjuntos de ecuaciones establecen que:
\begin{itemize}
 \item La producción total de cada industria (en términos físicos) es igual a la suma total de la cantidad de sus productos consumidos por las demás industrias.
 \item Bajo condiciones de equilibrio estacionario, el valor de los productos de cada industria es igual al valor de todos los bienes y servicios absorbidos por la misma industria.
 \item Describen la relación técnica entre la producción física de cada industria y la entrada de todos los diferentes elementos de costos absorbidos en la producción, es decir, intenta aproximar la función de producción industrial, véase \textcite{Leontief1937}.\footnote{Una descripción detallada del desarrollo histórico del modelo insumo--producto de Leontief y sus antecedentes puede encontrarse en el apéndice C de \textcite{Miller2012}. Una revisión por otras aproximaciones teóricas al tema de los flujos intersectoriales, específicamente la \textit{Tableau Économique} de Quesnay, los esquemas de reproducción de Marx y las ecuaciones de equilibrio de Walras, puede encontrarse en el capítulo 1 de \textcite{Mariña1993}.}
\end{itemize}

El modelo básico de insumo--producto considera un área económica particular, como un país, la cual es dividida en un conjunto de sectores económicos o industrias. Sus datos principales son los flujos intersectoriales en un periodo de tiempo específico, usualmente un año. Estos flujos son medidos en términos monetarios, debido a las dificultades de cuantificar flujos físicos cuando un sector vende diferentes bienes o servicios \parencite[p.~10--11]{Miller2012}.

La estructura matemática del sistema de insumo--producto básico está formada por un conjunto de $n$ ecuaciones lineales y $n$ variables, donde $n$ es el número de sectores económicos en el país \parencite[p.~10]{Miller2012}. Dados dos sectores económicos $i$ y $j$, sea $z_{ij}$ el valor de los bienes y servicios que el sector $i$ vende al sector $j$. Sea $x_i$ la producción total del sector $i$. Además de las industrias, el modelo insumo--producto considera la demanda final de cada sector $y_i$, bajo el supuesto de que son sectores exógenos que demandan bienes y servicios, pero no los ofrecen. La distribución de la producción de cada sector se representa con el siguiente sistema de ecuaciones:
\begin{eqnarray}
 x_1 &=& z_{11} + \ldots + z_{1j} + \ldots + z_{1n} + y_1\\
 \vdots & & \\
 x_i &=& z_{i1} + \ldots + z_{ij} + \ldots + z_{in} + y_i\\
 \vdots & & \\
 x_n &=& z_{n1} + \ldots + z_{nj} + \ldots + z_{nn} + y_n
\end{eqnarray}
Este sistema puede representarse de forma matricial como:
\begin{equation}
 \textbf{x} = \textbf{Z} \textbf{1} +  \textbf{y}
\end{equation}
donde 
\begin{eqnarray}
 \textbf{x} = \left[\begin{matrix} x_1\\ \vdots\\ x_n \end{matrix} \right], \quad \textbf{Z}=\left[\begin{matrix}
  z_{11} & \ldots & z_{1n}\\
  \vdots & \ddots & \vdots\\
  z_{n1} & \ldots & z_{nn}
 \end{matrix} \right],
 \qquad
 \textbf{y} = \left[\begin{matrix}
  y_1\\ \vdots\\ y_n
 \end{matrix} \right]
\end{eqnarray}
y $\textbf{1}$ representa un $n$-vector columna de unos \parencite[p.~11--12]{Miller2012}.

La demanda final suele dividirse en compras con fines de consumo o consumo en hogares $\textbf{C}$, consumo con fines de inversión o consumo privado $\textbf{I}$, consumo gubernamental $\textbf{G}$ y ventas en el extranjero o exportaciones $\textbf{E}$. Usando esta notación, la demanda final es igual al vector:
\begin{equation}
 \textbf{y} = \textbf{C}+\textbf{I}+\textbf{G}+\textbf{E}
\end{equation}

Para realizar su producción, cada sector $j$ compra bienes y servicios de los demás sectores, lo cual es representado por el vector columna 
\begin{equation}
 \textbf{z}_j = \left[\begin{matrix}
  z_{1j}\\
  \vdots\\
  z_{nj}
 \end{matrix} \right]
\end{equation}
este vector corresponde a los insumos intersectoriales o insumos intermedios, incluyendo los insumos provenientes del mismo sector $z_{jj}$. Además, el sector $j$ debe realizar pagos a otros elementos necesarios para la producción, como mano de obra, capital (pagos de intereses), servicios gubernamentales (impuestos), tierras (pagos de alquiler), emprendimientos (beneficios), los cuales son agrupados bajo el término de valor agregado y representados por $v_j$. Los bienes importados por el sector $j$ son representados por $m_j$. El valor agregado y las importaciones son denominados en conjunto como el sector de pagos \parencite[p.~12--14]{Miller2012}. Bajo esta notación, los desembolsos totales de un sector $j$ son:
\begin{equation}
 x_j = \sum_{i=1}^n z_{ij} + v_j + m_j
\end{equation}

Este modelo puede representarse mediante las tablas o matrices insumo--producto que consisten en una tabla de doble entrada, donde las columnas representan la composición del valor bruto de la producción desde la perspectiva del ingreso, y por filas se presenta su composición desde la perspectiva del gasto. Véase la figura \ref{Fig_MIP}. La \gls{mip} se divide en una matriz de insumos intermedios $\textbf{Z}$, un vector de demanda final $\textbf{y}$ y un vector del valor agregado $\textbf{VA}$ \parencite[p.~15]{Durán2021}. 

\begin{figure}[H]
 \centering
 \includegraphics[width=.7\linewidth]{MIP.png}
 \caption{Estructura básica de una \gls{mip} simétrica nacional.}
 \label{Fig_MIP}
\end{figure}

%El modelo insumo--producto es una de las fuentes de información estadística más importantes para estudiar la estructura económica de un país, ya que registra las transacciones que realizan entre sí los sectores económicos resultado de sus demandas y ofertas recíprocas de insumos\parencite{Mariñ11993}
% La elaboración estadística de las MIP se realiza de forma articulada con las cuentas nacionales y existe coherencia entre ambos sistemas. Algunos componentes del sistema de cuentas nacionales coinciden directamente con la MIP \parencite{Mariñ11993} (p-44--46)

\section{Matrices Insumo--Producto Multirregionales}

En las matrices insumo--producto nacionales el comercio internacional es considerado una variable exógena y no se incluye dentro de las relaciones intersectoriales. Las matrices insumo--producto multirregionales o multipaís (\acrshort{mriot})\footnote{Algunos autores usan el término GMRIO (\textit{Global Multiregional Input--Output}) para referirse a estos modelos, como en \textcite{Tukker2013}.} utiliza el marco teórico del modelo de Leontief, pero integran el comercio internacional de bienes y servicios intermedios, reflejando las relaciones intersectoriales entre los países o regiones.  Las fuentes de datos principales de las \gls{mip} multirregionales son las siguientes: a) Las estadísticas de las cuentas nacionales de cada país (cuadros de oferta y utilización o \glspl{mip} nacionales), b) datos sobre el comercio bilateral y, c) datos o supuestos sobre el origen del comercio internacional de los productos intermedios \parencite[p.~17]{Durán2021}.

La figura \ref{Fig_RIOT} muestra de forma simplificada la estructura de una \gls{mriot}. Una \gls{mriot} se compone de una matriz de demanda intermedia $\textbf{Z}_{m+1, m+1}$, donde $m$ es el número de países considerado. Cada elemento de la diagonal de $\textbf{Z}$ es, a su vez, una matriz $\textbf{Z}^{UU}$ de los insumos intermedios domésticos de un país $U$. Los elementos de $\textbf{Z}$ fuera de la diagonal son submatrices $n \times n$ de la forma:
\begin{equation}
 \textbf{Z}^{UV}=\left[\begin{matrix}
  z_{11}^{UV} & \ldots & z_{1n}^{UV} \\
  \vdots & \ddots & \vdots\\
  z_{n1}^{UV}  & \ldots & z_{nn}^{UV} 
 \end{matrix} \right],
\end{equation}
donde cada elemento $z_{ij}^{UV}$ representa los insumos exportados por el sector $i$ del país $U$ e importados por el sector $j$ del país $V$. Estas submatrices representan el flujo de valor en la cadena de valor global \parencite[p.~10]{Xing2022}. Este modelo incluye el comercio entre cada país considerado y todas las economías externas a la región, agrupadas bajo el término \gls{row}.

\begin{figure}[h]
 \centering
 \includegraphics[width=1\linewidth]{MRIOT.png}
 \caption{Estructura básica de una \gls{mriot} de 3 países.}
 \label{Fig_RIOT}
\end{figure}

Para cada país existe un vector fila de valor agregado $\textbf{VA}$ y un vector de producción bruta $\textbf{VBP}$. La demanda final se estructura en una matriz 
\begin{equation}\textbf{Y} = \left[ \textbf{Y}^{UV} \right]\end{equation}
conformada por submatrices $\textbf{Y}^{UV}$ con la información de la producción de cada sector de un país $U$ consumida en el país $V$ \parencite[p.~10]{Xing2022}.

Existe un tipo de matriz intermedia, denominada \gls{riot}. La \gls{riot} generalmente se concentra en una región del mundo y no incluye información sobre el consumo intermedio que es exportado por la región al resto del mundo, el consumo intermedio extrarregional ni la demanda final que es importada de fuera de la región. Se trata de una \gls{mriot} incompleta, que no está cerrada con el mundo \parencite[p.~19]{Durán2021}.

Las principales matrices insumo--producto multirregionales son:
\begin{itemize}
 \item \textit{Inter--Country Input--Output Tables} (\acrshort{icio}). Es una \gls{mriot} elaborada por la \gls{ocde}. Su última versión de 2025 muestra el flujo de bienes y servicios finales e intermedios de 80 países\footnote{Incluyen los 38 países de la  \gls{ocde}, y todos los países dentro del G20 y la \gls{ue}.} más el resto del mundo y 50 sectores, para el periodo 1995--2022. Los sectores se clasifican de acuerdo con la revisión 4 de la \gls{isic}. La versión extendida divide las economías de China y México en dos (mercado nacional y exportador).  Además del consumo intermedio, las tablas incluyen seis categorías de demanda final por país,\footnote{Gasto en consumo final de hogares; instituciones sin fines de lucro al servicio de los hogares; consumo final del gobierno general; formación bruta de capital fijo; cambios en inventarios y objetos de valor; y compras directas en el exterior por parte de residentes.} el valor agregado, y los impuestos menos subsidios a productos intermedios y finales por cada par país--industria \parencite{OECD-ICIO}.
 
 Existen varias versiones anteriores de esta base de datos. La versión de 1995 cubre 10 países de la \gls{ocde} y 35 industrias (\gls{isic} Rev. 2) durante el periodo de 1968 a 1990. La segunda versión de 2002 aumentó el número de países a 20 y de sectores a 42, con datos de 1992 a 1997. Mientras que la versión de 2006 incluía 37 países, 28 sectores y se amplió el periodo de análisis al año 2000 \parencite{Yamano2006}. En ambos casos la clasificación corresponde a la revisión 3 de la \gls{isic}. La edición de 2021 incluía datos de 71 países y 45 industrias para el periodo 1995--2018; mientras que la versión 2023 agregó 4 países más.
 
 Para su construcción, la \gls{ocde} utiliza tablas insumo--producto, y en su defecto, tablas de oferta y utilización generadas por los institutos nacionales de estadística. Los datos son armonizados en cuatro sentidos: homogeneizar los datos provenientes de diferentes países en la clasificación industrial de la \gls{isic}; estructurar los datos como una matriz industria--industria; expresar los valores en precios básicos;\footnote{Un precio básico es la cantidad de dinero que cobrará el productor al comprador por cada unidad de un bien o servicio producido menos los impuestos por pagar más las subvenciones por cobrar por el productor. Esto no incluye los gastos del transporte que el productor facture por separado \parencite[p.~8]{CEPAL2013}. No debe confundirse el término con los precios constantes a un año base.} y ajustarse con las estimaciones de los Sistemas de Cuentas Nacionales de cada país. Para más detalles metodológicos del proceso de armonización, véase \textcite{Yamano2006}.
 
 \item \textit{World Input--Output Database} (\acrshort{wiod}). Esta base de datos contiene los valores de transacciones anuales intersectoriales entre 43 países\footnote{28 países de la \gls{ue}, otros 15 países importantes del mundo y estimaciones sobre el resto del mundo en conjunto} y 56 sectores de acuerdo con \gls{isic} revisión 4, para el periodo 2000--2014, así como los flujos de cada industria--país a hogares, gobierno y usuarios de bienes de capital. Los valores de las tablas están expresados en precios de producción (o básicos) corrientes, en millones de dólares estadounidenses \parencite{Timmer2015}. %Uno of the objectives of WIOD is to permit the histórica comparison of the production, the agregated value and the commercial flujos in the economic environment }.
 Las tablas de insumo--producto mundial (\textit{World Input--Output Tables}, \acrshort{wiot}), núcleo de la WIOD, se construyeron a partir de series de tiempo de los cuadros de oferta y utilización (\textit{Supply and Use Tables}, \acrshort{sut}) nacionales, datos de las cuentas nacionales y estadísticas de comercio internacional. Los cuadros de oferta y utilización nacionales son utilizados como base para la construcción de una \acrshort{sut} mundial, la cual es transformada posteriormente en una \acrshort{wiot} \parencite{Dietzenbacher2013}. 
 
 Existe una versión de 2013, la cual cubre 35 industrias y 40 países para el periodo 1995--2011 y una versión de largo plazo que contempla el periodo de 1965--2000 con 25 países y 23 sectores económicos. En ambos la clasificación corresponde a la \gls{isic} revisión 3. El proyecto fue desarrollado por un equipo de la Universidad de Groningen, en Países Bajos y fue financiado por la \gls{ec} y la \acrlong{nwo}, \acrshort{nwo} \parencite{WIOD}.
 
 \item \textit{The Global Trade Analysis  Project} (\acrshort{gtap}) es una base de datos de paga que contiene tablas de insumo--producto nacionales, comercio bilateral y datos energéticos y ecológicos, de 141 países y 65 sectores, utilizando 5 años de referencia (2004, 2007, 2011, 2014, 2017) en su última versión de 2023 \parencite{GTAP11}. En total hay 11 versiones de la base de datos, que cubren el periodo entre 1993 y 2023, las cuales fueron agregando progresivamente países y sectores \parencite{GTAP_History}. Aunque la base \acrshort{gtap} no incluye directamente la matriz del comercio entre todos los países, puede construirse directamente una \gls{mriot} sin necesidad de hacer cálculos de equilibrio, como muestran en \textcite{Koopman2010, Peters2011, Andrew2013}. Esta base de datos no fue pensada para ofrecer tablas IO, sino para facilitar la simulación económica, asegurando que los datos económicos sean consistentes \parencite{Peters2011, Wiedmann2011}. Debido a que la base \acrshort{gtap} solo se compara con las estadísticas comerciales, los datos de oferta y demanda a nivel sectorial para países individuales pueden tener grandes discrepancias con los datos de las cuentas nacionales correspondientes. Este proyecto fue desarrollado por el Centro de Análisis del Comercio de la Universidad de Purdue \parencite{Jones2014}.
 
 \item \textit{International Input--Output Tables} del \gls{ide-jetro}. Esta base de datos está compuesta de tres \gls{mip} multirregionales:
 \begin{itemize}
  \item \textit{Asian International Input--Output Table} (AIIOT), 2005. Es una \gls{mriot} de 9 países asiáticos y los Estados Unidos y 76 sectores económicos. Considera otros elementos como importaciones y exportaciones a otros países específicos y al resto del mundo, valor agregado, demanda final y discrepancias estadísticas. Existen versiones anteriores de esta tabla, para los años 1985, 1990, 1995 y 2000.
  \item \textit{Transnational Interregional Input--Output Table for China, Japan and Korea} (TIIOT-CJK), 2005. La tabla están conformadas por los flujos insumo--producto entre 20 regiones de China, Japón y Corea del Sur, junto con las economías nacionales de Taiwán, los Estados Unidos, Hong Kong, y la economía agregada de los países de la \acrfull{asean} y el resto del mundo. En cuanto a las industrias, la tabla incluye dos clasificaciones, con 10 y 15 sectores agregados, respectivamente. Además, se reportan otros elementos en la tabla, como subdivisiones del valor agregado y la demanda final, costo de fletes, impuestos a importaciones y discrepancias estadísticas.
  \item \textit{BRICs International Input--Output Table 2005}. Esta tabla incluye los flujos insumo--producto de cada país de los BRICS\footnote{Foro político y económico formado por Brasil, Rusia, India, China y Sudáfrica (hasta 2023).} (excepto Sudáfrica), Japón, los Estados Unidos y la economía agregada de los países de la \acrlong{ue}, así como las importaciones y exportaciones al resto del mundo. Presenta dos clasificaciones de industrias en 7 y 25 sectores agregados \parencite{Meng2013, IDE-JETRO}.
 \end{itemize}
 Estas bases de datos reportan precios corrientes del productor, en miles de dólares estadounidenses. Una de sus desventajas es que sus datos solo corresponden a años de referencia específicos, a precios corrientes, por lo cual no pueden utilizarse en análisis a lo largo del tiempo \parencite[p.~578]{Timmer2015}.
 
 \item \textit{The Eora Global Supply Chain Database}. Es una base de datos de paga,\footnote{La base fue creada y es mantenida principalmente por investigadores de la Escuela de Física de la Universidad de Sydney, Austrialia, bajo financiamiento del \gls{arc} \cite[p.~103]{Miller2022}.} que proporciona una serie temporal de tablas insumo--producto, a precios corrientes, con cuentas satélite sociales y ambientales correspondientes a 192 economías y entre 26 y 511 sectores dependiendo del país.  La diferencia entre sectores se debe a que intentan mantener la clasificación sectorial de cada país cuando sea posible. La base de datos contiene series históricas de 1990 a 2022. Dado que considera prácticamente a todos los países del mundo, no construyen un agregado para el ``resto del mundo''. Existe una versión simplificada, denominada \textit{EORA26}, donde se usa una clasificación común de 26 sectores para todos los países, y las tablas de oferta y utilización del modelo completo se transforman en tablas de oferta y utilización simétricas producto--producto e industria--industria. Este modelo simplificado es más fácil de usar, pero menos preciso \parencite{EORA}. 

 Para su construcción combinan información de cuadros de oferta y utilización, así como tablas industria--industria y mercancía--mercancía de distintos países. La obtención de la \gls{mip} multirregional se realiza mediante estimaciones y el uso de métodos de optimización para asegurar el equilibrio de las tablas, a la vez que se intenta evitar la transformación de los datos originales de cada país, en aras de la transparencia
 \parencite{Lenzen2012, Lenzen2013}.
 

El objetivo del equipo que desarrolló este proyecto fue alcanzar el mayor grado de detalle respecto a los países, sectores y número de años. Además de las tablas de oferta y consumo y de insumo--producto, la base de datos \textit{Eora} tiene extensiones que incluyen consumo de energía, emisiones de gases de efecto invernadero y otras variables ambientales \parencite{Wiedmann2011}. Aunque sus datos cubren prácticamente a todas las economías del mundo, dependen de métodos de imputación para llenar los datos faltantes en países con sistemas estadísticos poco desarrollados \parencite[p.~579]{Timmer2015}.% Las tablas son presentadas en precios básicos.
 
 \item \textit{Asian Development Bank Multiregional Input--Output Tables} (\acrshort{adb}--MRIO) es una \gls{mip} multirregional elaborada por el \gls{adb} a partir de los datos de \gls{wiod}. Incluye 62 economías, más el resto del mundo, con desagregación por 35 industrias y 5 categorías de demanda final a precios constantes (base 2010). Estas tablas cubren el periodo de 2007 a 2023 y el año 2000. El \acrshort{adb}--MRIO también incluye tablas a precios corrientes para los mismos años, y para el periodo de 2017 a 2023 además incluye tablas con 72 economías \parencite{ABD}.
 
 \item \textit{Externality Data and Input--Output Tools for Policy Analysis} (EXIPOL). Es una \gls{riot} europea que cubre 27 países de la Unión Europea y otros 16 países (más el resto del mundo), así como 129 sectores industriales y productos. El objetivo principal del proyecto EXIPOL es mejorar la comprensión de los costos externos de las presiones ambientales. La tabla incluye información sobre el país de origen de cada producto y el sector receptor del país importador \parencite{Tukker2009, Wiedmann2011, Tukker2013}.
 
 \item Tablas regionales de insumo--producto de la zona euro:
 \begin{itemize}
  \item \textit{European Multi Regional Input Output}. Es una matriz insumo--producto regional de la Unión Europea, que incluye los flujos comerciales entre 272 regiones europeas (\acrshort{nuts}-2)\footnote{La \gls{nuts} es una clasificación jerárquica de 3 niveles, elaborada por la \gls{eurostat} para dividir geográficamente los países de la Unión Europea en regiones y subregiones para fines estadísticos y presupuestales. La \gls{nuts}-2 corresponde a regiones que tienen un rango poblacional entre 800 mil y 3 millones de habitantes \parencite{NUTS}.} y 10 sectores para el periodo 2008--2018 \parencite{Huang2023}.
  \item \textit{Eurostat's FIGARO\footnote{\textit{Full International and Global Accounts for Research in Input-Output Analysis}.} input--output tables}. Es una base de datos creada por la \gls{eurostat}, que incluye tablas de insumo--producto regional para 45 países (27 de la \gls{ue}), más el resto del mundo, y 64 industrias, para los años 2010--2022, así como componentes de la demanda final y el valor agregado \parencite{FIGARO}.
  
  %\item EUREGIO: The Construction of a Global IO Database With Regional Detail for Europe for 2000–2010. https://papers.ssrn.com/sol3/papers.cfm?abstract_id=3285818
 \end{itemize}
\end{itemize}

En la tabla \ref{tab_MRIOT} se comparan las características de las últimas versiones de cada \gls{mriot}. Se muestra el nombre de la base de datos y el año o número de su última versión. La tercera columna indica el número de países cubierto, mientras que la cuarta columna representa el número de sectores o industrias que incluye la base de datos. La quinta columna contiene el número de componentes de la demanda final que son considerados en la \gls{mriot}, mientras que la última columna muestra los años disponibles.


\begin{table}[h]
\centering
\caption{Características de las principales matrices insumo--producto multirregionales.}
\label{tab_MRIOT}
\begin{tabular}{@{} p{4cm} c c c p{2cm} p{2cm} @{}} 
\hline
\textbf{Base de datos} & \multicolumn{1}{l}{\textbf{Versión}} & \multicolumn{1}{l}{\textbf{Economías}} & \multicolumn{1}{l}{\textbf{Sectores}} & \textbf{Demanda final} & \textbf{Periodo} \\ \hline
OECD--ICIO & 2025 & 81 & 50 & 6 & 1995--2022 \\
WIOD & 2016 & 43 & 56 & 5 & 2000--2014 \\
GTAP & v11 & 141 & 65 & -- & 2004, 2007, 2011, 2014, 2017 \\
AIIOT & 2005 & 10 & 76 & 5 & 2005 \\
IDE--JETRO (BRICS) & 2005 & 7 & 10, 15 & 5 & 2005 \\
TIIOT-CJK & 2005 & 7 & 10, 15 & 5 & 2005 \\
EORA & v199.82 & 190 & 26--511 & 6 & 1990--2022 \\
ABD--MRIO & 2023 & 62 & 35 & 5 & 2000, 2007--2023 \\
EXIPOL & -- & 43 & 129 & -- & -- \\
European MRIO &  2023 & 27 & 10 &  5 &  2008--2018\\
FIGARO & -- & 45 & 64 & 4 & 2010--2022 \\ \hline
\end{tabular}
\end{table}
% TIIOT-CJK - 20 regiones
% EUROPEAN MRIO 272 regiones


La necesidad de armonizar y consolidar datos de diferentes sistemas estadísticos nacionales ha dificultado que los institutos nacionales de estadística y organismos supranacionales contribuyan en la construcción de \gls{mip} multirregionales, con la excepción de la base de datos de la \gls{ocde} \parencite[p.~7]{Tukker2013}. 

En \textcite{Jones2014} comparan tres \gls{mip} multirregionales (\gls{wiod}--2013, \acrshort{icio} versión 2013 y \acrshort{gtap} v8) con los datos de estadísticas oficiales para el periodo 2004--2009. Después de armonizar las tres tablas en la misma clasificación de países y sectores, verifican su coherencia con estadísticas oficiales de las \gls{unna}: \gls{pib} de cada país (basado en cuenta de gastos), demanda interna final y balanza de pagos. Los autores encontraron que los valores del \gls{pib} y de la demanda final calculados a partir de los datos de las tablas \gls{icio} son consistentes con los datos informados por la UNNA, presentando pequeñas desviaciones. Por su parte, los cálculos de ambos indicadores basados en \gls{wiod} fluctúan alrededor de los valores de referencia de las \gls{unna} ($\pm 1\%$ para el 70\% de las economías consideradas). Mientras que la base \gls{gtap} presenta valores consistentemente inferiores a los informados por las \gls{unna}, con diferencias significativas para algunos países. Las diferencias entre las estadísticas oficiales y las tres MIP multirregionales son mucho mayores y fluctúan más en los valores comerciales (importaciones y exportaciones) que en términos del \gls{pib} o la demanda final. Las tablas \gls{icio} son las que presentan las menores diferencias con los datos oficiales: para la mayoría de las economías no pertenecientes a la \gls{ue}, presentan diferencias negativas pequeñas o insignificantes, mientras que las mayores diferencias se presentan en países miembros de la \gls{ue}. Los datos basados en \gls{wiod} son consistentemente menores que las estadísticas oficiales, particularmente en países de la \gls{ue}. Los autores sugieren que estas desviaciones pueden deberse, en gran parte, al registro sobre reexportaciones y comercio de tránsito que tiene lugar dentro de la \gls{ue}.

En esta tesis se decidió trabajar con las tablas \gls{icio} versión 2025 por las siguientes razones. En primer lugar, por su amplia cobertura en términos de países, sectores y años. En segundo lugar, por su metodología armonizada con los sistemas de cuentas nacionales y las matrices insumo--producto oficiales de cada país, así como con la clasificación internacional de la \gls{isic} revisión 4. Además, sus datos son consistentes con las estimaciones del \gls{pib}, la demanda final y la balanza de pagos de los países. %de acuerdo con el trabajo de \parencite{Jones2014}. 
Finalmente, es una base de datos abierta y gratuita, elaborada por un organismo supranacional.

En la tabla \ref{Tab_Cobertura_PorContinente} se presenta el número de países considerados  en la base de datos \gls{icio}--\gls{ocde} versión 2025, agrupados por continente. Por su parte, la figura \ref{Fig_Cobertura_ICIO25} muestra en rojo los países cubiertos por esta base de datos, mientras que los países no resaltados son agrupados bajo la categoría ``resto del mundo''. Las tablas \ref{ICIO25_Paises} y \ref{ICIO25_Industrias} del anexo \ref{anexo1} muestran el listado detallado de los países y sectores incluidos en la base. 

\begin{table}[H]
	\centering
	\caption{Cobertura por continente de la base ICIO--OCDE versión 2025}
	\label{Tab_Cobertura_PorContinente}
	\begin{tabular}{lr}
		\toprule
		\textbf{Continente} & \textbf{Cantidad} \\ \midrule
		Europa              & 33                                    \\
		Asia                & 25                                    \\
		África              & 11                                    \\
		Sudamérica          & 5                                     \\
		América del Norte   & 4                                     \\
		Oceanía             & 2                                     \\ \bottomrule
	\end{tabular}	
\end{table}

\begin{figure}[H]
 \centering
 \includegraphics[width=\linewidth]{Cobertura2025.png}
 \caption{Cobertura geográfica de la base ICIO--OCDE versión 2025.}
 \label{Fig_Cobertura_ICIO25}
\end{figure}



%\textcolor{red}{[Agregar mapa de países considerados y tabla de sectores.]}




